%----------------------------------------------------------------------------------------
% BASIS EN DIMENSIE. Alle opgaven zijn eigen werk.
%
% Opbouw: Probleem, Start, leerstof, en op een nieuwe pagina de Opdrachten, Koppel en Test jezelf.
% Labels van dit hoofdstuk beginnen met h2-.
%----------------------------------------------------------------------------------------

\mychapter{2}{Basis en dimensie}

\begin{probleem}
In een \textit{magisch vierkant} van 3 bij 3 hebben de drie rijen, de drie kolommen en de twee diagonalen dezelfde som. Het oudste bekende komt uit China, met som 15. Ook het vierkant met overal een 1 is magisch, met som 3. Tel je twee magische vierkanten vakje per vakje op, dan krijg je weer een magisch vierkant, en dat geldt ook voor een veelvoud. We laten alle reële getallen toe, ook negatieve en breuken.

\figuren{\vierkantdrie{4 & 9 & 2 \\ \hline 3 & 5 & 7 \\ \hline 8 & 1 & 6}
\quad $+\ \ 2\,\cdot$\quad
\vierkantdrie{1 & 1 & 1 \\ \hline 1 & 1 & 1 \\ \hline 1 & 1 & 1}
\quad $=$\quad
\vierkantdrie{6 & 11 & 4 \\ \hline 5 & 7 & 9 \\ \hline 10 & 3 & 8}}

\noindent Hoeveel magische vierkanten moet je kennen om er met optellen en vermenigvuldigen met getallen \textit{alle} magische vierkanten mee te maken? En hoeveel vakjes moet je invullen voordat de rest vastligt?
\end{probleem}

\lkr{Leerplan Complementair (III-WiCo-d): LPD 33, met alle wenken: lineair onafhankelijk, voortbrengend, basis, de dimensie (met het bewijs dat ze niet van de basis afhangt), een oneindigdimensionale vectorruimte, en coördinaten ten opzichte van een basis. Voorkennis: hoofdstuk 1, en uit pakket 4 de rang, de determinant en de stelling dat een vierkante matrix met rang $n$ inverteerbaar is en een determinant verschillend van 0 heeft (Matrices, hoofdstuk 3). De criteria met de rang en de determinant, en de dimensie van de oplossingen van $AX = O$, verbinden deze cursus met Matrices.}

In hoofdstuk 1 zag je dat een paar vectoren een hele vectorruimte kunnen voortbrengen. Hoe weinig vectoren zijn er nodig? Is een vector overbodig, dan kun je hem weglaten. Blijven er alleen nodige vectoren over, dan heb je een \textit{basis}, en dan kun je elke vector op precies één manier als combinatie schrijven. De coëfficiënten zijn dan coördinaten, zoals in een assenstelsel.

Het aantal vectoren in een basis is een eigenschap van de vectorruimte zelf: het is altijd hetzelfde, welke basis je ook kiest. Dat getal heet de \textit{dimensie}. Het vlak heeft dimensie 2 en de ruimte dimensie 3, zoals je verwacht. Maar ook de magische vierkanten hebben een dimensie, en die beantwoordt de vraag van het probleem.

\begin{oefeningen}[Start]
\oef{h2-start-1}

\oef{h2-verken}
\end{oefeningen}

\lkr{Laat oefening \ref{oef:h2-verken} maken vóór de uitleg. Met twee vectoren in het vlak is de combinatie uniek, met drie niet: dan kun je ook de nulvector maken met coëfficiënten die niet nul zijn. Dat is het idee achter lineair (on)afhankelijk.}

%========================================================================================
\section{Lineair onafhankelijk}

\subsection{Overbodige vectoren}

In opdracht \ref{oef:h1-o6} van hoofdstuk 1 brachten $(1, -1, 2)$ en $(-2, 2, -4)$ samen maar een rechte voort: de tweede is $-2$ keer de eerste, en brengt niets nieuws. Hoe zie je in het algemeen of er een vector overbodig is? Een vector is overbodig als hij een combinatie is van de andere. Dat kun je ook anders zeggen: dan is er een combinatie van alle vectoren die de nulvector geeft, zonder dat alle coëfficiënten nul zijn. Uit $\vec v_2 = -2\vec v_1$ volgt zo $2\vec v_1 + \vec v_2 = \nulv$. Die tweede vorm is eerlijker: ze kiest geen vector uit.

\begin{theorie}
De vectoren $\vec v_1, \ldots, \vec v_k$ zijn \textbf{\textcolor{accent}{lineair onafhankelijk}} als
\[ \formule{c_1\vec v_1 + c_2\vec v_2 + \cdots + c_k\vec v_k = \nulv \ \Rightarrow \ c_1 = c_2 = \cdots = c_k = 0} \]
Anders zijn ze \textbf{\textcolor{accent}{lineair afhankelijk}}: dan is er een combinatie gelijk aan $\nulv$ waarin niet alle coëfficiënten nul zijn.

\begin{stelling}{Afhankelijke vectoren}
De vectoren $\vec v_1, \ldots, \vec v_k$ (met $k \geqslant 2$) zijn lineair afhankelijk als en slechts als één van hen een lineaire combinatie is van de andere.
\end{stelling}

\begin{bewijs}
Stel $c_1\vec v_1 + \cdots + c_k\vec v_k = \nulv$ met bijvoorbeeld $c_1 \neq 0$. Breng de andere termen naar rechts en deel door $c_1$:
\bewijsregel{\vec v_1 = -\frac{c_2}{c_1}\vec v_2 - \cdots - \frac{c_k}{c_1}\vec v_k .}
Stel omgekeerd dat $\vec v_1 = d_2\vec v_2 + \cdots + d_k\vec v_k$. Dan is
\bewijsregel{(-1)\vec v_1 + d_2\vec v_2 + \cdots + d_k\vec v_k = \nulv ,}
een combinatie met een coëfficiënt $-1 \neq 0$. Voor een andere vector dan $\vec v_1$ gaat het net zo.
\end{bewijs}
\end{theorie}

Kort: afhankelijk betekent dat er een vector overbodig is, onafhankelijk dat elke vector nodig is. Een paar gevallen zie je meteen. Eén vector is onafhankelijk als hij niet de nulvector is. Twee vectoren zijn afhankelijk als ze collineair zijn. Drie vectoren in de ruimte zijn afhankelijk als ze coplanair zijn. En een stel met de nulvector erin is altijd afhankelijk, want $1 \cdot \nulv = \nulv$.

\figuren{\begin{tikzpicture}[scale=.6,font=\scriptsize]
  \draw[->,lijngrijs] (-.3,0) -- (3.3,0);
  \draw[->,lijngrijs] (0,-.3) -- (0,2.4);
  \vpijl{(0,0)}{(1.2,.8)}
  \vpijl[fel!80!black]{(0,0)}{(2.7,1.8)}
  \node[above left] at (1.2,.8) {$\vec v_1$};
  \node[above] at (2.7,1.8) {$\vec v_2$};
  \node[font=\footnotesize] at (1.5,-.8) {collineair: afhankelijk};
\end{tikzpicture}
\qquad
\begin{tikzpicture}[scale=.6,font=\scriptsize]
  \draw[->,lijngrijs] (-.3,0) -- (3.3,0);
  \draw[->,lijngrijs] (0,-.3) -- (0,2.4);
  \vpijl{(0,0)}{(2.4,.6)}
  \vpijl[fel!80!black]{(0,0)}{(.7,1.9)}
  \node[below] at (2.4,.6) {$\vec v_1$};
  \node[left] at (.7,1.9) {$\vec v_2$};
  \node[font=\footnotesize] at (1.5,-.8) {onafhankelijk};
\end{tikzpicture}
\qquad
\begin{tikzpicture}[ruimte,scale=.75,font=\scriptsize]
  \fill[accent!14] (2.2,-1,-.69) -- (2.2,2.3,.135) -- (-1.4,2.3,.855) -- (-1.4,-1,.03) -- cycle;
  \assendrie{2.4}
  \vpijl{(0,0,0)}{(1.5,0,-.3)}
  \vpijl{(0,0,0)}{(0,1.6,.4)}
  \vpijl[fel!80!black]{(0,0,0)}{(1.2,1.4,.11)}
  \node[font=\footnotesize] at (0,1.1,-1.85) {coplanair: afhankelijk};
\end{tikzpicture}}

\voorbeeld{}{Zijn $(1, 2, 1)$, $(2, 1, 0)$ en $(0, 3, 2)$ lineair onafhankelijk?\label{vb:h2-onafh}}{%
\uitleg{We zoeken alle combinaties die de nulvector geven. Dat is een homogeen stelsel met de coëfficiënten als onbekenden. Is de enige oplossing $c_1 = c_2 = c_3 = 0$, dan zijn de vectoren onafhankelijk.}
\begin{stappenbreed}
\begin{cases} c_1 + 2c_2 = 0 \\ 2c_1 + c_2 + 3c_3 = 0 \\ c_1 + 2c_3 = 0 \end{cases} & $c_1(1, 2, 1) + c_2(2, 1, 0) + c_3(0, 3, 2) = (0, 0, 0)$, per coördinaat. \\[5ex]
c_1 = -2c_3, \quad c_2 = c_3 & Uit de derde en dan de eerste vergelijking. \\[1ex]
-4c_3 + c_3 + 3c_3 = 0 & De tweede vergelijking geldt voor elke $c_3$. \\[1ex]
-2\,(1, 2, 1) + (2, 1, 0) + (0, 3, 2) = (0, 0, 0) & Neem $c_3 = 1$: een combinatie die niet alleen nullen heeft.
\end{stappenbreed}
De vectoren zijn lineair afhankelijk. Zo is $(0, 3, 2) = 2\,(1, 2, 1) - (2, 1, 0)$: de drie vectoren liggen in één vlak door de oorsprong.}

\nujij{Zijn $(1, 0, 2)$ en $(3, 0, 5)$ lineair onafhankelijk?}{Ja: ze zijn niet collineair, want $3 \cdot 2 = 6 \neq 5$. Uit $c_1 + 3c_2 = 0$ en $2c_1 + 5c_2 = 0$ volgt $c_1 = c_2 = 0$.}

\begin{oefeningen}
\oef{h2-onafh}
\end{oefeningen}

\subsection{Onafhankelijk en voortbrengend met de rang}

In $\R^n$ werk je met kolommen. Zet de vectoren $\vec v_1, \ldots, \vec v_k$ als kolommen in een $n \times k$-matrix $A$. Dan is de combinatie $c_1\vec v_1 + \cdots + c_k\vec v_k$ precies het product $AC$, met $C$ de kolom van de coëfficiënten. Zo wordt de vraag of de vectoren onafhankelijk zijn, de vraag hoeveel oplossingen $AC = O$ heeft. En dat vertelt de rang (Matrices, hoofdstuk 3).

\begin{theorie}
\begin{stelling}{Onafhankelijk en rang}
De vectoren $\vec v_1, \ldots, \vec v_k$ van $\R^n$ zijn lineair onafhankelijk als en slechts als de matrix $A$ met die vectoren als kolommen rang $\formule{\rang A = k}$ heeft.
\end{stelling}

\begin{bewijs}
De vectoren zijn onafhankelijk als en slechts als het homogene stelsel $AC = O$ alleen de oplossing $C = O$ heeft. Een homogeen stelsel is nooit strijdig, en het heeft $k - \rang A$ parameters. Er is dus alleen de oplossing $C = O$ als en slechts als $k - \rang A = 0$.
\end{bewijs}
\end{theorie}

Omdat de rang hoogstens het aantal rijen is, volgt meteen: meer dan $n$ vectoren in $\R^n$ zijn altijd afhankelijk. Drie vectoren in het vlak zijn dus afhankelijk, zoals in de Verken.

\voorbeeld{}{Zijn $(1, 0, 2, 1)$, $(0, 1, 1, -1)$ en $(1, 1, 0, 1)$ lineair onafhankelijk in $\R^4$?\label{vb:h2-rang}}{%
\uitleg{We zetten de vectoren als kolommen in een $4 \times 3$-matrix en bepalen de rang. Is die 3, dan zijn ze onafhankelijk.}
\begin{stappenbreed}
\km{1 & 0 & 1 \\ 0 & 1 & 1 \\ 2 & 1 & 0 \\ 1 & -1 & 1} \to \km{1 & 0 & 1 \\ 0 & 1 & 1 \\ 0 & 1 & -2 \\ 0 & -1 & 0} & $R_3 - 2R_1$ en $R_4 - R_1$. \\[5ex]
\to \km{1 & 0 & 1 \\ 0 & 1 & 1 \\ 0 & 0 & -3 \\ 0 & 0 & 1} & $R_3 - R_2$ en $R_4 + R_2$: in elke kolom een spil. \\[5ex]
\rang A = 3 & Drie spillen, en er zijn drie vectoren.
\end{stappenbreed}
De vectoren zijn lineair onafhankelijk.}

Met dezelfde matrix kun je ook nagaan of vectoren $\R^n$ voortbrengen. Dan moet het stelsel $AC = B$ een oplossing hebben voor elke kolom $B$.

\begin{theorie}
\begin{stelling}{Voortbrengend en rang}
De vectoren $\vec v_1, \ldots, \vec v_k$ brengen $\R^n$ voort als en slechts als de matrix $A$ met die vectoren als kolommen rang $\formule{\rang A = n}$ heeft.
\end{stelling}

\begin{bewijs}
Een vector $B$ van $\R^n$ is een combinatie van de kolommen als en slechts als $AC = B$ niet strijdig is.

Is $\rang A = n$, dan heeft ook $[A \,|\, B]$ rang $n$: ze heeft $n$ rijen, en minstens de $n$ spillen van $A$. Het stelsel is dus niet strijdig, voor elke $B$.

Is $\rang A < n$, rijherleid dan $A$ tot $R$: de onderste rij van $R$ is een nulrij. Zet er een kolom $E$ naast met onderaan een 1 en verder nullen. Het stelsel $RC = E$ is strijdig, want de laatste vergelijking zegt $0 = 1$. Voer nu de rijbewerkingen in omgekeerde volgorde ongedaan op $[R \,|\, E]$: links krijg je $A$ terug, rechts een kolom $B$. Rijbewerkingen veranderen de oplossingen niet, dus is $AC = B$ ook strijdig, en $B$ is geen combinatie.
\end{bewijs}
\end{theorie}

Omdat de rang hoogstens het aantal kolommen is, hebben minder dan $n$ vectoren nooit genoeg om $\R^n$ voort te brengen. Twee vectoren brengen de ruimte niet voort: ze vullen hoogstens een vlak.

\nujij{Brengen $(1, 2, 0)$, $(0, 1, 1)$ en $(1, 3, 1)$ de ruimte $\R^3$ voort?}{Nee. De derde is de som van de eerste twee, dus na rijherleiden is er een nulrij: de rang is 2.}

\begin{oefeningen}
\oef{h2-rang}
\end{oefeningen}

\subsection{Functies}

In een vectorruimte van functies zijn er geen coördinaten om in een matrix te zetten. Toch kun je onafhankelijkheid nagaan: een combinatie is de nulfunctie als ze 0 is \textit{in elk getal} $x$. Vul je een paar handige getallen in, dan krijg je een stelsel voor de coëfficiënten.

\voorbeeld{}{Toon aan dat de functies $\e^x$ en $\e^{2x}$ lineair onafhankelijk zijn.\label{vb:h2-functies}}{%
\uitleg{Stel dat $a\,\e^x + b\,\e^{2x} = 0$ voor elke $x$. Dan geldt het zeker voor twee gekozen getallen. We kiezen $x = 0$ en $x = \ln 2$, zodat de machten eenvoudig worden.}
\begin{stappen}
a + b = 0 & $x = 0$: $\e^0 = 1$. \\[1ex]
2a + 4b = 0 & $x = \ln 2$: $\e^{\ln 2} = 2$ en $\e^{2\ln 2} = 4$. \\[1ex]
b = 0, \quad a = 0 & Trek twee keer de eerste af van de tweede: $2b = 0$.
\end{stappen}
Alleen $a = b = 0$ geeft de nulfunctie: ze zijn onafhankelijk.}

Andersom bewijs je dat functies afhankelijk zijn met een gelijkheid die voor elke $x$ geldt, zoals een formule uit de goniometrie. Zo zijn $1$, $\sin^2 x$ en $\cos^2 x$ afhankelijk, want $\sin^2 x + \cos^2 x - 1 = 0$ voor elke $x$. Getallen invullen helpt daar niet: die tonen alleen dat de coëfficiënten aan een paar voorwaarden voldoen.

\begin{werkwijze}
Schrijf $c_1\vec v_1 + \cdots + c_k\vec v_k = \nulv$ uit. In $\R^n$: zet de vectoren als kolommen in een matrix en bepaal de rang; is die $k$, dan zijn ze onafhankelijk. Bij functies: vul een paar getallen in. Geeft dat alleen $c_1 = \cdots = c_k = 0$, dan zijn ze onafhankelijk. Om afhankelijkheid te tonen, geef je één combinatie die $\nulv$ is en niet alleen nullen heeft.
\end{werkwijze}

\begin{oefeningen}
\oef{h2-functies}
\end{oefeningen}

%========================================================================================
\section{Basis en coördinaten}

\subsection{Een basis}

Een stel vectoren dat voortbrengend is, laat je elke vector maken. Een stel dat onafhankelijk is, heeft geen overbodige vectoren. Heb je allebei, dan kun je elke vector maken met zo weinig mogelijk vectoren.

\begin{theorie}
Een \textbf{\textcolor{accent}{basis}} van een vectorruimte $V$ is een rij vectoren $B = (\vec b_1, \ldots, \vec b_n)$ die voortbrengend is voor $V$ en lineair onafhankelijk.

\begin{stelling}{Coördinaten}
$(\vec b_1, \ldots, \vec b_n)$ is een basis van $V$ als en slechts als elke vector $\vec v$ van $V$ op precies één manier te schrijven is als
\[ \formule{\vec v = c_1\vec b_1 + c_2\vec b_2 + \cdots + c_n\vec b_n} \]
\end{stelling}

\begin{bewijs}
Stel dat $B$ een basis is. Omdat $B$ voortbrengend is, is elke $\vec v$ een combinatie. Stel dat er twee zijn: $\vec v = c_1\vec b_1 + \cdots + c_n\vec b_n = d_1\vec b_1 + \cdots + d_n\vec b_n$. Trek af:
\bewijsregel{(c_1 - d_1)\vec b_1 + \cdots + (c_n - d_n)\vec b_n = \nulv .}
Omdat $B$ onafhankelijk is, zijn alle $c_i - d_i = 0$: de twee combinaties zijn dezelfde.

Omgekeerd: is elke vector een combinatie, dan is $B$ voortbrengend. Omdat ook $\nulv$ maar op één manier een combinatie is, namelijk met alle coëfficiënten 0, is $B$ onafhankelijk.
\end{bewijs}

De getallen $c_1, \ldots, c_n$ zijn de \textbf{\textcolor{accent}{coördinaten}} van $\vec v$ ten opzichte van $B$. Je noteert \quad $\formule{\coord{\vec v}{B} = (c_1, \ldots, c_n)}$
\end{theorie}

Een basis is een rij: de volgorde telt, want ze bepaalt welke coördinaat bij welke vector hoort. Elke vectorruimte uit hoofdstuk 1 heeft een basis waarin de coördinaten de getallen zijn die je al kende. Dat is de \textit{standaardbasis}.

\begin{center}
\arrayrulecolor{grijs}\renewcommand{\arraystretch}{1.45}
\begin{tabular}{|c|c|c|}
\hline
\rowcolor{accent!14} vectorruimte & standaardbasis & een vector en zijn coördinaten \\ \hline
$\R^3$ & $(1, 0, 0),\ (0, 1, 0),\ (0, 0, 1)$ & $(4, -1, 2)$ \ heeft coördinaten \ $(4, -1, 2)$ \\ \hline
$\Rx{2}$ & $1,\ x,\ x^2$ & $5 - 3x + x^2$ \ heeft coördinaten \ $(5, -3, 1)$ \\ \hline
$\Rmn{2}{2}$ & $\km{1 & 0 \\ 0 & 0}, \km{0 & 1 \\ 0 & 0}, \km{0 & 0 \\ 1 & 0}, \km{0 & 0 \\ 0 & 1}$ & $\km{2 & 7 \\ 0 & -1}$ \ heeft coördinaten \ $(2, 7, 0, -1)$ \\ \hline
\end{tabular}
\end{center}

Bij een andere basis horen andere coördinaten. In de figuur is $B = \bigl((2, 1), (1, 2)\bigr)$. Het scheve rooster bestaat uit de lijnen $c_1 \vec b_1 + t\,\vec b_2$ en $t\,\vec b_1 + c_2\vec b_2$ met gehele $c_1$ en $c_2$, zoals het gewone rooster bij de standaardbasis. De vector $(5, 4)$ ligt op het kruispunt van de lijnen met $c_1 = 2$ en $c_2 = 1$.

\voorbeeld{}{Bepaal de coördinaten van $\vec v = (5, 4)$ ten opzichte van $B = \bigl((2, 1), (1, 2)\bigr)$.\label{vb:h2-coord}}{%
\uitleg{We zoeken $c_1$ en $c_2$ met $c_1(2, 1) + c_2(1, 2) = (5, 4)$: een stelsel met twee vergelijkingen.}
\naast{\begin{stappenvrij}{4.2cm}{4.8cm}
\begin{cases} 2c_1 + c_2 = 5 \\ c_1 + 2c_2 = 4 \end{cases} & Per coördinaat. \\[3ex]
3c_2 = 3 & Twee keer de tweede min de eerste. \\[1ex]
c_2 = 1, \quad c_1 = 2 & Invullen in de tweede.
\end{stappenvrij}
Dus $\coord{\vec v}{B} = (2, 1)$: $(5, 4) = 2\,(2, 1) + (1, 2)$. Ten opzichte van de standaardbasis zijn de coördinaten $(5, 4)$.}{\begin{tikzpicture}[scale=.5,font=\scriptsize]
  \begin{scope}\clip (-.6,-.6) rectangle (6.4,5.6);\scheefrooster{2}{1}{1}{2}{-2}{4}{-2}{4}\end{scope}
  \roosterleeg{-.6}{1}{6.4}{-.6}{1}{5.6}
  \draw[grijs,dashed] (4,2) -- (5,4);
  \vpijl{(0,0)}{(2,1)}
  \vpijl{(0,0)}{(1,2)}
  \vpijl[fel!80!black]{(0,0)}{(5,4)}
  \node[below right] at (2,1) {$\vec b_1$};
  \node[above left] at (1,2) {$\vec b_2$};
  \node[above] at (5,4) {$(5, 4)$};
  \foreach \t in {2,4,6}{\node[below,inner sep=1.5pt] at (\t,0) {$\t$};}
  \foreach \t in {2,4}{\node[left,inner sep=1.5pt] at (0,\t) {$\t$};}
\end{tikzpicture}}}

In een andere vectorruimte gaat het net zo. Een basis van veeltermen die zich rond een ander punt dan 0 schikt, is handig als je een veelterm rond dat punt wilt bekijken.

\voorbeeld{}{Bepaal de coördinaten van $p(x) = 2x^2 - 3x + 4$ ten opzichte van $B = \bigl(1,\ x - 1,\ (x - 1)^2\bigr)$.\label{vb:h2-coordveelterm}}{%
\uitleg{We zoeken $a$, $b$ en $c$ met $a + b(x - 1) + c(x - 1)^2 = 2x^2 - 3x + 4$. We werken het linkerlid uit en vergelijken de coëfficiënten, te beginnen bij de hoogste macht.}
\begin{stappenbreed}
a + bx - b + cx^2 - 2cx + c = cx^2 + (b - 2c)x + (a - b + c) & Werk uit en orden. \\[1ex]
c = 2 & De coëfficiënt van $x^2$. \\[1ex]
b - 4 = -3 \ \Rightarrow \ b = 1 & De coëfficiënt van $x$. \\[1ex]
a - 1 + 2 = 4 \ \Rightarrow \ a = 3 & De constante term.
\end{stappenbreed}
Dus $\coord{p}{B} = (3, 1, 2)$: $p(x) = 3 + (x - 1) + 2(x - 1)^2$. Controle: $p(1) = 2 - 3 + 4 = 3$, en dat is de eerste coördinaat.}

\begin{werkwijze}
Schrijf $\vec v = c_1\vec b_1 + \cdots + c_n\vec b_n$ met onbekende coëfficiënten. Stel het gelijk per coördinaat (of per macht van $x$), en los het stelsel op. Controleer door de combinatie uit te rekenen.
\end{werkwijze}

\begin{oefeningen}
\oef{h2-coord}
\end{oefeningen}

\subsection{Een basis van \texorpdfstring{$\R^n$}{Rn} herkennen}

Hoe zie je snel of $n$ vectoren van $\R^n$ een basis vormen? Dan is de matrix met die vectoren als kolommen vierkant, en dan helpt de determinant.

\begin{theorie}
\begin{stelling}{Basis van $\R^n$}
De vectoren $\vec v_1, \ldots, \vec v_n$ vormen een basis van $\R^n$ als en slechts als de $n \times n$-matrix $A$ met die vectoren als kolommen een determinant heeft met $\formule{\det A \neq 0}$
\end{stelling}

\begin{bewijs}
De vectoren vormen een basis als en slechts als elke $B$ van $\R^n$ op precies één manier een combinatie is (stelling over coördinaten), dus als $AC = B$ bepaald is voor elke $B$. Volgens de stelling over rang, determinant en inverse uit Matrices is dat zo als en slechts als $\det A \neq 0$.
\end{bewijs}
\end{theorie}

\voorbeeld{}{Vormen $(1, 1, 0)$, $(1, 0, 1)$ en $(0, 1, 1)$ een basis van $\R^3$?\label{vb:h2-basisdet}}{%
\uitleg{Drie vectoren in $\R^3$: de matrix is vierkant, dus berekenen we de determinant met de regel van Sarrus.}
\begin{stappenbreed}
\det\km{1 & 1 & 0 \\ 1 & 0 & 1 \\ 0 & 1 & 1} = (0 + 0 + 0) - (0 + 1 + 1) = -2 & De vectoren staan als kolommen. \\[2ex]
-2 \neq 0 & Dus een basis.
\end{stappenbreed}
Ze vormen een basis van $\R^3$. Omdat $\det A\tr = \det A$, mag je de vectoren ook als rijen zetten.}

\begin{oefeningen}
\oef{h2-basisdet}
\end{oefeningen}

%========================================================================================
\section{Dimensie}

\subsection{Elke basis is even groot}

Het vlak $\R^2$ heeft de basis $\bigl((1, 0), (0, 1)\bigr)$, maar ook $\bigl((2, 1), (1, 2)\bigr)$ en oneindig veel andere. Telkens met twee vectoren. Kan een basis van $\R^2$ ook uit drie vectoren bestaan, of uit één? In $\R^n$ zegt de rang: nee. Maar hoe zit het in een vectorruimte zonder coördinaten, zoals de magische vierkanten? De sleutel is de volgende stelling: wie door $m$ vectoren voortgebracht wordt, heeft geen plaats voor meer dan $m$ onafhankelijke vectoren.

\begin{theorie}
\begin{stelling}{Te veel vectoren}
Wordt $V$ voortgebracht door $m$ vectoren, dan zijn $k > m$ vectoren van $V$ altijd lineair afhankelijk.
\end{stelling}

\begin{bewijs}
Stel $V = \vb{\vec v_1, \ldots, \vec v_m}$ en neem $\vec w_1, \ldots, \vec w_k$ in $V$, met $k > m$. Elke $\vec w_j$ is een combinatie van de $\vec v_i$: noem
\bewijsregel{\vec w_j = a_{1j}\vec v_1 + a_{2j}\vec v_2 + \cdots + a_{mj}\vec v_m ,}
en $A = [a_{ij}]$ de $m \times k$-matrix van die coëfficiënten. Het homogene stelsel $AC = O$ heeft $m$ vergelijkingen en $k$ onbekenden. Zijn rang is hoogstens $m < k$, dus heeft het minstens één parameter: er is een oplossing $C = (c_1, \ldots, c_k) \neq O$. Dan is
\bewijsregel{c_1\vec w_1 + \cdots + c_k\vec w_k = \sum_{j=1}^{k} c_j \sum_{i=1}^{m} a_{ij}\vec v_i = \sum_{i=1}^{m} \Bigl(\sum_{j=1}^{k} a_{ij} c_j\Bigr)\vec v_i = \sum_{i=1}^{m} 0 \cdot \vec v_i = \nulv ,}
want de som tussen haakjes is rij $i$ van $AC = O$. Dat is een combinatie gelijk aan $\nulv$ met niet alle $c_j = 0$.
\end{bewijs}
\end{theorie}

Het idee van het bewijs: we vertalen de vraag naar een homogeen stelsel met meer onbekenden dan vergelijkingen, en zo'n stelsel heeft altijd een oplossing die niet nul is. Uit de stelling volgt dat het aantal vectoren in een basis vastligt.

\begin{theorie}
\begin{stelling}{Even grote bases}
Heeft $V$ een basis met $n$ vectoren, dan heeft elke basis van $V$ precies $n$ vectoren.
\end{stelling}

\begin{bewijs}
Neem een basis met $n$ vectoren en een andere met $k$ vectoren. De eerste brengt $V$ voort, en de vectoren van de tweede zijn onafhankelijk. Volgens de vorige stelling is dus $k \leqslant n$. Wissel de rollen om: dan is $n \leqslant k$. Dus $k = n$.
\end{bewijs}

Heeft $V$ een basis met $n$ vectoren, dan is $n$ de \textbf{\textcolor{accent}{dimensie}} van $V$: \quad $\formule{\dim V = n}$

\vspace{1mm}
De vectorruimte $\{\nulv\}$ heeft dimensie 0. Een vectorruimte die door een eindig aantal vectoren wordt voortgebracht, heet \textbf{\textcolor{accent}{eindigdimensionaal}}; anders heet ze \textbf{\textcolor{accent}{oneindigdimensionaal}}.
\end{theorie}

De dimensie telt dus hoeveel getallen je nodig hebt om een vector vast te leggen: zijn coördinaten. Met de standaardbases vind je meteen de dimensie van de bekende vectorruimten.

\begin{center}
\arrayrulecolor{grijs}\renewcommand{\arraystretch}{1.45}
\begin{tabular}{|c|c|c|}
\hline
\rowcolor{accent!14} vectorruimte & dimensie & voorbeeld \\ \hline
$\R^n$ & $n$ & $\dim \R^3 = 3$ \\ \hline
$\Rx{n}$ & $n + 1$ & $\dim \Rx{2} = 3$: \ $1, x, x^2$ \\ \hline
$\Rmn{m}{n}$ & $mn$ & $\dim \Rmn{2}{2} = 4$ \\ \hline
\end{tabular}
\end{center}

\voorbeeld{}{Bepaal een basis en de dimensie van $W = \bigl\{p \in \Rx{2} \bigm| p(1) = 0\bigr\}$.\label{vb:h2-dimdeel}}{%
\uitleg{We schrijven een willekeurige veelterm van $W$ op met zo weinig mogelijk vrije getallen. Die vrije getallen worden de coördinaten.}
\begin{stappenbreed}
p(x) = a + bx + cx^2, \quad a + b + c = 0 & De voorwaarde $p(1) = 0$. \\[1ex]
p(x) = (-b - c) + bx + cx^2 = b(x - 1) + c(x^2 - 1) & Vul $a = -b - c$ in, en zet $b$ en $c$ samen. \\[1ex]
W = \vb{x - 1,\ x^2 - 1} & Voortbrengend. \\[1ex]
b(x - 1) + c(x^2 - 1) = 0 \ \Rightarrow \ c = 0,\ b = 0 & Onafhankelijk: kijk naar de coëfficiënt van $x^2$, dan die van $x$.
\end{stappenbreed}
$(x - 1, x^2 - 1)$ is een basis van $W$, en $\dim W = 2$.}

\begin{oefeningen}
\oef{h2-dimensie}
\end{oefeningen}

\subsection{Oneindigdimensionaal}

Heeft elke vectorruimte een eindige dimensie? Nee. Neem $\R[x]$, alle veeltermen, van elke graad. Elke veelterm is een combinatie van $1, x, x^2, \ldots$, maar daar heb je er oneindig veel van nodig.

\begin{theorie}
\begin{stelling}{Alle veeltermen}
De vectorruimte $\R[x]$ van alle veeltermen is oneindigdimensionaal.
\end{stelling}

\begin{bewijs}
Stel dat een eindig aantal veeltermen $p_1, \ldots, p_m$ de ruimte $\R[x]$ voortbrengt. Noem $d$ de grootste graad van die veeltermen. Een combinatie $c_1p_1 + \cdots + c_mp_m$ heeft dan graad hoogstens $d$. Maar $x^{d+1}$ ligt in $\R[x]$ en heeft graad $d + 1$: het is geen combinatie. Dat is een contradictie.
\end{bewijs}
\end{theorie}

Ook de functies van $\R$ naar $\R$ en de rijen vormen oneindigdimensionale vectorruimten: ze bevatten de veeltermen, of de rijen $(1, 0, 0, \ldots)$, $(0, 1, 0, \ldots)$, \ldots, waarvan je er telkens zoveel als je wilt kunt nemen die onafhankelijk zijn. Volgens de stelling over te veel vectoren kan zo'n vectorruimte door geen enkel eindig aantal vectoren voortgebracht worden. Vectorruimten van functies zijn het onderwerp van de functionaalanalyse, die de natuurkunde gebruikt in de kwantummechanica.

\subsection{Een basis maken}

Hoe vind je een basis? Je kunt vertrekken van vectoren die voortbrengend zijn en er schrappen, of van onafhankelijke vectoren en er bijvoegen.

\begin{theorie}
\begin{stelling}{Uitdunnen}
Brengen $\vec v_1, \ldots, \vec v_k$ de vectorruimte $V \neq \{\nulv\}$ voort, dan vormen sommige van hen een basis van $V$.
\end{stelling}

\begin{bewijs}
Zijn de vectoren afhankelijk, dan is er een, zeg $\vec v_k$, een combinatie van de andere. Elke combinatie van alle $k$ vectoren is dan ook een combinatie van de eerste $k - 1$: vervang $\vec v_k$ door zijn combinatie. Dus $\vb{\vec v_1, \ldots, \vec v_{k-1}} = V$. Herhaal dat zolang de vectoren afhankelijk zijn. Na hoogstens $k - 1$ stappen blijven er onafhankelijke vectoren over die $V$ voortbrengen: een basis.
\end{bewijs}
\end{theorie}

Een eindigdimensionale vectorruimte heeft dus altijd een basis, en dus een dimensie. Om de overbodige vectoren in $\R^n$ te vinden, gebruik je de rijgereduceerde vorm. Zet de vectoren als kolommen in $A$ en rijherleid tot $R$. Rijbewerkingen veranderen de oplossingen van $AC = O$ niet: de kolommen van $A$ hebben dus dezelfde lineaire verbanden als die van $R$. In $R$ zie je ze meteen. De kolommen met een spil zijn onafhankelijk, en elke andere kolom is een combinatie van de spilkolommen ervoor.

\voorbeeld{}{Bepaal een basis en de dimensie van $W = \vb{(1, 2, 1), (2, 1, 0), (0, 3, 2), (1, 0, 1)}$.\label{vb:h2-uitdunnen}}{%
\uitleg{We zetten de vier vectoren als kolommen en rijherleiden. De spilkolommen van $A$ vormen een basis.}
\begin{stappenbreed}
\km{1 & 2 & 0 & 1 \\ 2 & 1 & 3 & 0 \\ 1 & 0 & 2 & 1} \to \km{1 & 0 & 2 & 0 \\ 0 & 1 & -1 & 0 \\ 0 & 0 & 0 & 1} & Gauss-Jordan. Spillen in de kolommen 1, 2 en 4. \\[4ex]
(0, 3, 2) = 2\,(1, 2, 1) - (2, 1, 0) & Kolom 3 van $R$ is $2 \cdot$ kolom 1 $-$ kolom 2: dat verband geldt ook in $A$. \\[1ex]
\bigl((1, 2, 1), (2, 1, 0), (1, 0, 1)\bigr) & De spilkolommen van $A$: een basis van $W$.
\end{stappenbreed}
$\dim W = 3$, en dus is $W = \R^3$ (zie verder). De derde vector was overbodig: dat zagen we al in voorbeeld \ref{vb:h2-onafh}.}

\begin{oefeningen}
\oef{h2-uitdunnen}
\end{oefeningen}

Omgekeerd kun je een onafhankelijk stel aanvullen. Voeg je een vector toe die nog niet voortgebracht wordt, dan blijft het stel onafhankelijk.

\begin{theorie}
\begin{stelling}{Aanvullen}
Zijn $\vec v_1, \ldots, \vec v_k$ lineair onafhankelijk en ligt $\vec v$ niet in $\vb{\vec v_1, \ldots, \vec v_k}$, dan zijn $\vec v_1, \ldots, \vec v_k, \vec v$ lineair onafhankelijk.
\end{stelling}

\begin{bewijs}
Stel $c_1\vec v_1 + \cdots + c_k\vec v_k + c\,\vec v = \nulv$. Was $c \neq 0$, dan was $\vec v = -\frac{c_1}{c}\vec v_1 - \cdots - \frac{c_k}{c}\vec v_k$ een combinatie, en dat is niet zo. Dus $c = 0$, en dan is $c_1\vec v_1 + \cdots + c_k\vec v_k = \nulv$, zodat ook alle $c_i = 0$.
\end{bewijs}
\end{theorie}

\subsection{Een basis herkennen}

Ken je de dimensie, dan hoef je voor een basis maar de helft na te gaan: het juiste aantal onafhankelijke vectoren brengt vanzelf voort, en het juiste aantal voortbrengende vectoren is vanzelf onafhankelijk.

\begin{theorie}
\begin{stelling}{$n$ onafhankelijke vectoren}
Is $\dim V = n$, dan vormen $n$ lineair onafhankelijke vectoren van $V$ een basis.
\end{stelling}

\begin{bewijs}
Stel dat $\vec v_1, \ldots, \vec v_n$ onafhankelijk zijn, maar $V$ niet voortbrengen. Neem dan $\vec v \in V$ buiten $\vb{\vec v_1, \ldots, \vec v_n}$. Door de stelling over aanvullen zijn $\vec v_1, \ldots, \vec v_n, \vec v$ onafhankelijk: $n + 1$ onafhankelijke vectoren in een vectorruimte die door een basis van $n$ vectoren wordt voortgebracht. Dat kan niet, door de stelling over te veel vectoren.
\end{bewijs}
\end{theorie}

\begin{theorie}
\begin{stelling}{$n$ voortbrengende vectoren}
Is $\dim V = n$, dan vormen $n$ vectoren die $V$ voortbrengen een basis.
\end{stelling}

\begin{bewijs}
Volgens de stelling over uitdunnen vormen sommige van de $n$ vectoren een basis. Elke basis heeft $n$ vectoren, dus je hebt er geen geschrapt: de $n$ vectoren zelf vormen de basis.
\end{bewijs}
\end{theorie}

Zo zijn $1 + x$, $x + x^2$ en $1 + x^2$ een basis van $\Rx{2}$ zodra je weet dat ze onafhankelijk zijn: het zijn er drie, en $\dim \Rx{2} = 3$.

\subsection{Deelruimten en hun dimensie}

Een deelruimte van $\R^3$ is een rechte door de oorsprong, een vlak door de oorsprong, \ldots\ Zijn er nog andere? De dimensie geeft het antwoord.

\begin{theorie}
\begin{stelling}{Dimensie van een deelruimte}
Is $W$ een deelruimte van $V$ en $\dim V = n$, dan is $W$ eindigdimensionaal en $\formule{\dim W \leqslant n}$. Is $\dim W = n$, dan is $W = V$.
\end{stelling}

\begin{bewijs}
Is $W = \{\nulv\}$, dan is $\dim W = 0$. Neem anders $\vec w_1 \neq \nulv$ in $W$. Brengt die $W$ niet voort, neem dan $\vec w_2 \in W$ buiten $\vb{\vec w_1}$, en zo verder. Door de stelling over aanvullen zijn $\vec w_1, \ldots, \vec w_k$ telkens onafhankelijk, ook in $V$. Volgens de stelling over te veel vectoren is dus $k \leqslant n$: na hoogstens $n$ stappen brengen ze $W$ voort, en vormen ze een basis van $W$. Dus $\dim W \leqslant n$.

Is $\dim W = n$, dan is een basis van $W$ een stel van $n$ onafhankelijke vectoren in $V$, en dus een basis van $V$. Dan is $W = \vb{\vec w_1, \ldots, \vec w_n} = V$.
\end{bewijs}
\end{theorie}

\naast{\begin{minipage}{10.2cm}
Zo vind je alle deelruimten van $\R^3$. Een deelruimte heeft dimensie 0, 1, 2 of 3. Dimensie 0 is $\{\nulv\}$. Dimensie 1 is $\vb{\vec v}$ met $\vec v \neq \nulv$: een rechte door de oorsprong. Dimensie 2 is $\vb{\vec v_1, \vec v_2}$ met twee onafhankelijke vectoren: een vlak door de oorsprong. En dimensie 3 is $\R^3$ zelf. Er zijn dus geen andere. Net zo heeft $\R^2$ alleen $\{\nulv\}$, de rechten door de oorsprong en $\R^2$ zelf als deelruimten, zoals beloofd in hoofdstuk 1.
\end{minipage}}{\begin{tikzpicture}[ruimte,scale=.62,font=\scriptsize]
  \fill[accent!14] (2.2,-1.4,0) -- (2.2,2.4,0) -- (-1.4,2.4,0) -- (-1.4,-1.4,0) -- cycle;
  \assendrie{2.6}
  \draw[fel!80!black,very thick] (-1.2,-1.1,-1.3) -- (1.5,1.4,1.65);
  \fill[black] (0,0,0) circle (2pt);
  \node[font=\footnotesize] at (1.6,-1.6,0) {dim 2};
  \node[font=\footnotesize,right] at (1.5,1.4,1.65) {dim 1};
  \node[font=\footnotesize,below right] at (0,0,0) {dim 0};
\end{tikzpicture}}

\subsection{De oplossingen van een homogeen stelsel}

In hoofdstuk 1 zag je dat de oplossingen van $AX = O$ een deelruimte vormen. Welke dimensie? Uit Matrices weet je dat er $n - \rang A$ parameters zijn. Elke parameter levert een vector van een basis.

\begin{theorie}
\begin{stelling}{Dimensie van de oplossingen}
Is $A$ een $m \times n$-matrix, dan vormen de oplossingen van $AX = O$ een deelruimte van $\R^n$ met dimensie $\formule{n - \rang A}$
\end{stelling}

\begin{bewijs}
Noem $p = n - \rang A$. Na rijherleiden geef je de $p$ onbekenden zonder spil de waarden $t_1, \ldots, t_p$, en volgen de andere onbekenden daaruit, lineair. Noem $X_j$ de oplossing met $t_j = 1$ en de andere parameters 0. Dan is elke oplossing
\bewijsregel{X = t_1X_1 + t_2X_2 + \cdots + t_pX_p ,}
dus $X_1, \ldots, X_p$ brengen de oplossingen voort. Ze zijn ook onafhankelijk: op de plaats van de onbekende met parameter $t_j$ heeft $X_j$ een 1 en heeft elke andere $X_i$ een 0. Is $t_1X_1 + \cdots + t_pX_p = O$, dan staat op die plaats $t_j = 0$, voor elke $j$.
\end{bewijs}
\end{theorie}

\voorbeeld{}{Bepaal een basis en de dimensie van de oplossingen van $\begin{cases} x_1 + 2x_2 - x_3 + x_4 = 0 \\ 2x_1 + 4x_2 + x_3 + 5x_4 = 0 \end{cases}$\label{vb:h2-oplossingen}}{%
\uitleg{We rijherleiden, kiezen de onbekenden zonder spil als parameters, en splitsen de oplossing per parameter.}
\begin{stappenbreed}
\km{1 & 2 & -1 & 1 \\ 2 & 4 & 1 & 5} \to \km{1 & 2 & 0 & 2 \\ 0 & 0 & 1 & 1} & $R_2 - 2R_1$, dan $\frac13 R_2$ en $R_1 + R_2$. Spillen bij $x_1$ en $x_3$. \\[2ex]
x_2 = s, \ x_4 = t: \quad x_1 = -2s - 2t, \ x_3 = -t & Twee parameters: $4 - 2 = 2$. \\[1ex]
X = s\,(-2, 1, 0, 0) + t\,(-2, 0, -1, 1) & Splits per parameter.
\end{stappenbreed}
Een basis is $\bigl((-2, 1, 0, 0), (-2, 0, -1, 1)\bigr)$, en de dimensie is 2. Controle met de eerste vector in beide vergelijkingen: $-2 + 2 = 0$ en $-4 + 4 = 0$.}

\begin{werkwijze}
Een basis van een deelruimte die met vergelijkingen gegeven is: los het homogene stelsel op met parameters, en splits de oplossing per parameter. De vectoren bij de parameters vormen een basis, en hun aantal is de dimensie. Is de deelruimte gegeven met voortbrengende vectoren: dun ze uit met de spilkolommen.
\end{werkwijze}

\begin{oefeningen}
\oef{h2-oplossingen}
\end{oefeningen}

\subsection{Terug naar het probleem}

\voorbeeld{}{Terug naar het probleem. In een \textit{magisch vierkant} van 3 bij 3 hebben de drie rijen, de drie kolommen en de twee diagonalen dezelfde som. Het oudste bekende komt uit China, met de rijen $4, 9, 2$; $3, 5, 7$ en $8, 1, 6$ en som 15. Ook het vierkant met overal een 1 is magisch. Tel je twee magische vierkanten vakje per vakje op, dan krijg je weer een magisch vierkant, en dat geldt ook voor een veelvoud. We laten alle reële getallen toe. Hoeveel magische vierkanten moet je kennen om er met optellen en vermenigvuldigen met getallen \textit{alle} magische vierkanten mee te maken? En hoeveel vakjes moet je invullen voordat de rest vastligt?\label{vb:h2-terug}}{%
\uitleg{De magische vierkanten zijn $3 \times 3$-matrices. Zeggen dat alle acht sommen gelijk zijn, is zeggen dat zeven verschillen nul zijn, zoals (rij 1) $-$ (rij 2) $= 0$: dat zijn zeven homogene vergelijkingen in de negen getallen. De magische vierkanten vormen dus een deelruimte $M$ van $\Rmn{3}{3}$, met dimensie hoogstens 9. We zoeken een basis van $M$. Noem de vakjes zoals in de figuur, en $s$ de magische som. We tonen eerst dat het middelste vakje een derde van $s$ is. Daarna drukken we alle vakjes uit in $a$, $b$ en $e$, en controleren we dat zo'n vierkant echt magisch is.}
\naast{\begin{minipage}{11.6cm}
\begin{stappenvrij}{5.4cm}{5.7cm}
4s = (a + \cdots + i) + 3e = 3s + 3e & Tel de middelste rij, de middelste kolom en de twee diagonalen op: elk vakje één keer, $e$ vier keer. Dus $s = 3e$. \\[1ex]
c = 3e - a - b, \quad i = 2e - a & Rij 1 en de diagonaal $a$, $e$, $i$. \\[1ex]
h = 2e - b, \quad g = a + b - e & Kolom 2 en de diagonaal $c$, $e$, $g$. \\[1ex]
d = 4e - 2a - b, \quad f = 2a + b - 2e & Kolom 1, dan rij 2.
\end{stappenvrij}
\end{minipage}}{\vierkantdrie{a & b & c \\ \hline d & e & f \\ \hline g & h & i}}
\vspace{1mm}
\noindent Controle van de twee sommen die we niet gebruikten: rij 3 geeft $(a + b - e) + (2e - b) + (2e - a) = 3e$ en kolom 3 geeft $(3e - a - b) + (2a + b - 2e) + (2e - a) = 3e$. Elk magisch vierkant is dus
\[ \vierkantdrie[20mm]{a & b & 3e - a - b \\ \hline 4e - 2a - b & e & 2a + b - 2e \\ \hline a + b - e & 2e - b & 2e - a} \]
of, gesplitst naar $a$, $b$ en $e$,
\[ a \cdot \vierkantdrie{1 & 0 & -1 \\ \hline -2 & 0 & 2 \\ \hline 1 & 0 & -1} \ + \ b \cdot \vierkantdrie{0 & 1 & -1 \\ \hline -1 & 0 & 1 \\ \hline 1 & -1 & 0} \ + \ e \cdot \vierkantdrie{0 & 0 & 3 \\ \hline 4 & 1 & -2 \\ \hline -1 & 2 & 2} \]
en elke zulke combinatie is magisch. Noem de drie vierkanten $A$, $B$ en $E$: ze brengen $M$ voort. Ze zijn onafhankelijk, want in de combinatie $aA + bB + eE$ staan $a$, $b$ en $e$ zelf in drie vakjes: is ze nul, dan zijn $a$, $b$ en $e$ nul. Dus is $(A, B, E)$ een basis van $M$, en $\dim M = 3$.

Je hebt dus drie magische vierkanten nodig, en drie vakjes ($a$, $b$ en $e$) leggen de rest vast. Ze zijn de coördinaten ten opzichte van $(A, B, E)$. Het vierkant uit China heeft coördinaten $(4, 9, 5)$: vul in, en je vindt $c = 2$, $d = 3$, $f = 7$, $g = 8$, $h = 1$ en $i = 6$.}

Niet elke drie vakjes leggen een magisch vierkant vast. Ken je $a$, $e$ en $i$, dan weet je niets extra, want $i = 2e - a$ volgt al uit $a$ en $e$. De vraag welke drie vakjes werken, is de vraag welke drie coördinaten onafhankelijk zijn: dat zie je in de opdrachten. Op de voorkant van dit boek staat een magisch vierkant van 4 bij 4, uit de prent \textit{Melencolia I} van Albrecht Dürer (1514). De magische vierkanten van 4 bij 4 vormen een vectorruimte van dimensie 8.

\begin{oefeningen}
\oef{h2-magisch}
\end{oefeningen}

%========================================================================================
\opdrachtenkop

\begin{oefeningen}
\oef{h2-o1}

\oef{h2-o2}

\oef{h2-o3}

\oef{h2-o4}

\needspace{14\baselineskip}
\oef{h2-o5}

\oef{h2-o6}

\oef{h2-o7}

\oef{h2-o8}

\oef{h2-o9}

\oef{h2-o10}

\oef{h2-o11}

\oef{h2-o12}

\oef{h2-o13}

\oef{h2-o14}

\oef{h2-o15}

\oef{h2-o16}

\oef{h2-o17}

\oef{h2-o18}

\oef{h2-o19}

\oef{h2-o20}

\oef{h2-o21}

\oef{h2-o22}
\end{oefeningen}

%----------------------------------------------------------------------------------------
% KOPPEL: verbanden met andere leerstof
%----------------------------------------------------------------------------------------
\needspace{20\baselineskip}
\opdrachtsectie{Koppel}

\lkr{De eerste oefening veronderstelt de lineaire differentiaalvergelijkingen van de tweede orde uit de cursus Differentiaalvergelijkingen (complementair gedeelte). Complexe getallen en Ruimtemeetkunde komen uit pakket 4.}

\begin{oefeningen}
\oef{h2-koppel-1}

\oef{h2-koppel-2}

\oef{h2-koppel-3}
\end{oefeningen}

%----------------------------------------------------------------------------------------
\needspace{30\baselineskip}
\opdrachtsectie{Test jezelf}

\begin{oefeningen}
\oef{h2-t1}

\oef{h2-t2}

\oef{h2-t3}

\oef{h2-t4}

\oef{h2-t5}

\oef{h2-t6}

\oef{h2-t7}
\end{oefeningen}

\herhaalschema{%
  \fv{oef:h2-t1}{vb:h2-onafh}{oef:h2-onafh}%
  \fv{oef:h2-t2}{vb:h2-functies}{oef:h2-functies}%
  \fv{oef:h2-t3}{vb:h2-basisdet}{oef:h2-basisdet}%
  \fv{oef:h2-t4}{vb:h2-coord}{oef:h2-coord}%
  \fv{oef:h2-t5}{vb:h2-coordveelterm}{oef:h2-coord}%
  \fv{oef:h2-t6}{vb:h2-oplossingen}{oef:h2-oplossingen}%
  \fv{oef:h2-t7}{vb:h2-dimdeel}{oef:h2-dimensie}}

\ifzelfstudie
\section*{Maak je eigen samenvatting}
Sluit dit boek en schrijf in je schrift: wat lineair onafhankelijk en afhankelijk betekenen en wat dat met overbodige vectoren te maken heeft, hoe je het nagaat in $\R^n$ (rang) en bij functies, wanneer vectoren $\R^n$ voortbrengen, wat een basis en coördinaten zijn en waarom coördinaten uniek zijn, de standaardbases, de determinant als test voor een basis van $\R^n$, waarom elke basis even groot is, de dimensie van $\R^n$, $\Rx{n}$ en $\Rmn{m}{n}$, een oneindigdimensionale vectorruimte, uitdunnen en aanvullen, de deelruimten van $\R^3$, en de dimensie van de oplossingen van $AX = O$. Vergelijk daarna met de kaders.
\fi
